\subsection{Pantallas coordinadas y dualidad de radios recíprocos}
\label{exc:pantallas-dualidad}

La reducción seleccionada por \(K\) y la construcción de Leech
proceden de lecturas distintas de la misma incidencia. Para precisar
su coordinación, se reúnen aquí los mapas algebraicos que conservan
ambos datos. El espacio de rango once y el retículo de rango veinticuatro
no se identifican; se especifican sus respectivos cocientes y la
involución que actúa en el factor de rutas y memoria.

\subsubsection{Residuo, cociente y cambio de sección}

La evaluación aritmética levantada de APP conserva
\(n=9Q+r\), con \(r\in\{1,\ldots,9\}\) en la sección
residual positiva. El transporte TPK conserva por separado estas
coordenadas, su orientación y la memoria de sus cambios. Para
representar un intercambio entre las dos coordenadas se requiere
un dominio estable por intercambio. Se usa por ello el ambiente
\(\ZZ^2\), cuyos enteros firmados prolongan las coordenadas de
la carta; ello no afirma que todas sus parejas sean trayectorias
admisibles de una misma historia TPK.

En \(\mathscr D_0=\ZZ^2\times\RR_{>0}\), la realización cuadrática
declarada es
\begin{equation}
 q_R(m,w)=\frac{m^2}{R^2}+w^2R^2,
 \qquad \mathsf T_0(m,w,R)=(w,m,R^{-1}).
 \label{eq:dualidad-normalizada-k}
\end{equation}
La variable \(R\) es una escala adimensional de la carta de lectura.
Los símbolos \(m,w\) designan aquí las dos coordenadas enteras
transportadas. Su reconocimiento posterior como momento y enrollamiento
mantiene esa distinción entre construcción y realización.

\begin{proposition}[Intercambio y forma cuadrática]
\label{prop:dualidad-radios-k}
La aplicación \(\mathsf T_0\) es una involución y satisface
\[
 q_{R^{-1}}(w,m)=q_R(m,w).
\]
En \(\ell^2(\ZZ^2)\), el operador unitario
\(U_T|m,w\rangle=|w,m\rangle\) entrelaza los operadores
diagonales \(H(R)|m,w\rangle=q_R(m,w)|m,w\rangle\):
\[
 U_TH(R)U_T^{-1}=H(R^{-1}).
\]
\end{proposition}
\begin{proof}
Dos intercambios recuperan las coordenadas y dos inversiones recuperan
\(R\). La sustitución da
\(q_{R^{-1}}(w,m)=w^2R^2+m^2/R^2\). El intercambio permuta
la base ortonormal, luego es unitario, y esa igualdad verifica el
entrelazamiento sobre cada vector de la base. Se extiende a los
dominios de los operadores diagonales por la misma igualdad de pesos.
\end{proof}

La sección residual importa. Para \(n=9\), las parejas
\((r,Q)=(9,0)\) y \((0,1)\) representan el mismo entero;
sus evaluaciones \(81/R^2\) y \(R^2\) no coinciden para todo
\(R\). El cambio de sección se incorpora exactamente mediante
\[
 L_9=\ZZ(9,-1),\qquad S(x_1,x_2)=(x_2,x_1).
\]
Definimos en las clases del cociente
\begin{equation}
 \mathcal E_R^{L_9}([x])=\min_{\ell\in L_9}q_R(x+\ell),
 \qquad
 \mathsf T_{\mathrm{cov}}([x]_{L_9},R)
       =([Sx]_{SL_9},R^{-1}).
 \label{eq:dualidad-cociente-k}
\end{equation}

\begin{proposition}[Dualidad compatible con el cambio de sección]
Las dos expresiones de \eqref{eq:dualidad-cociente-k} están bien
definidas. Las secciones positiva y cero, con el correspondiente
incremento de cociente, representan la misma clase. Además,
\[
 \mathcal E_{R^{-1}}^{SL_9}([Sx])=\mathcal E_R^{L_9}([x]).
\]
\end{proposition}
\begin{proof}
Cambiar \(x\) por \(x+\ell_0\) permuta los candidatos al mínimo.
Éste existe porque \(q_R(x+k(9,-1))\) es una cuadrática coerciva
en \(k\in\ZZ\). El cambio \((9,Q)\mapsto(0,Q+1)\)
tiene diferencia \((9,-1)\). Finalmente,
\[
 \min_{\ell\in L_9}q_{R^{-1}}(Sx+S\ell)
 =\min_{\ell\in L_9}q_R(x+\ell).
\]
Como \(S^2=I\), el intercambio inverso recupera la clase y la escala.
\end{proof}

La igualdad depende de transportar también la relación de equivalencia,
de \(L_9\) a \(SL_9\). Así se conserva la información del cambio
de sección que una mera sustitución del representante \(9\) por
\(0\) perdería~\cite{hmtedicionintegral}.

\subsubsection{La pantalla reticular y su coordinación con \(K\)}

Sea \(\Lambda=N_\alpha\), el retículo de Leech construido en
\ref{exc:leech}. Añadimos el plano hiperbólico integral
\(U=\ZZ e_+\oplus\ZZ e_-\), con
\(e_+^2=e_-^2=0\) y \(\langle e_+,e_-\rangle=1\).
Entonces \(L_{26}=\Lambda\oplus U\) tiene firma \((25,1)\) y
\[
 e_+^\perp=\Lambda\oplus\ZZ e_+,
 \qquad e_+^\perp/\ZZ e_+\cong\Lambda.
\]
En efecto, si \(y=\lambda+ae_++be_-\), la condición
\(\langle y,e_+\rangle=0\) equivale a \(b=0\), y el cociente
elimina exactamente el sumando \(ae_+\). Su mapa es
\(q_{24}(\lambda+ae_+)=\lambda\).

Para la reducción \(P_{10}\) ya demostrada, sea
\[
 \mathscr G_K=\{(v_{11},P_{10}v_{11}):v_{11}\in V_{11}\}.
\]
Definimos
\begin{equation}
 \mathcal R_K(v,d,y)
 =\bigl((P_{11}v,P_{10}v),d,q_{24}(y)\bigr),
 \label{eq:k-mapa-pantallas}
\end{equation}
con dominio \(\RR^{12}\times\mathscr D_0\times e_+^\perp\)
y codominio \(\mathscr G_K\times\mathscr D_0\times\Lambda\).

\begin{theorem}[Isomorfismo de las pantallas cocientadas]
\label{thm:k-pantallas-coordinadas}
El mapa \eqref{eq:k-mapa-pantallas} induce un isomorfismo
\[
 (\RR^{12}/\RR\mathbf1)\times\mathscr D_0
       \times(e_+^\perp/\ZZ e_+)
 \ \cong\ \mathscr G_K\times\mathscr D_0\times\Lambda.
\]
Entrelaza las transformaciones que actúan por \(\mathsf T_0\)
en el factor \(\mathscr D_0\) y por la identidad en los demás.
\end{theorem}
\begin{proof}
\(P_{11}\) induce un isomorfismo de \(\RR^{12}/\RR\mathbf1\)
sobre \(V_{11}\). Como \(P_{10}P_{11}=P_{10}\), la aplicación
\([v]\mapsto(P_{11}v,P_{10}v)\) identifica ese cociente con
el grafo \(\mathscr G_K\); la inversa toma la primera componente
y su clase. El cociente reticular acaba de identificarse mediante
\(q_{24}\). El factor restante se conserva por identidad.
El producto de las tres aplicaciones prueba el isomorfismo. La
involución y los proyectores actúan en factores distintos, por lo
que sus composiciones conmutan. La forma cuadrática se conserva
por la proposición \ref{prop:dualidad-radios-k}.
\end{proof}

El grafo retiene \(v_{11}\) junto con su proyección. Por ello este
isomorfismo no afirma que \(V_{11}\) y \(V_{10}\) sean isomorfos:
la reducción aislada tiene núcleo \(\ell_K\), mientras el grafo
conserva la coordenada previa. Esta precisión identifica qué conserva
cada mapa y permite enlazar las pantallas \(11\to10\) y
\(26\to24\) sin confundirlas.

En conjunto, \(K\) determina la dirección de reducción; la bandera
dodecafásica determina el vecino reticular utilizado para Moonshine;
y el intercambio de las coordenadas de ruta y memoria conserva la
forma cuadrática bajo radios recíprocos. La aplicación a una teoría
física de cuerdas o a un espacio de dimensión once requiere además
su acción, campos y operadores de realización. El resultado reunido
aquí es el mapa algebraico con sus pruebas, que constituye un
antecedente explícito para ese desarrollo posterior.
